\section{The \dataset Benchmark}
\label{sec:dataset}

\subsection{Operational Event Graph}

Given a context image $x^g$ and a local evidence crop $x^l$, the pair task predicts an event set $Y\subseteq\mathcal{Y}$. Each event $y\in Y$ is represented as
\begin{equation}
  y=(d,e,s,r,c),
\end{equation}
where $d$ is the operational domain, $e$ the entity or region, $s$ its state, $r$ a spatial or functional relation, and $c$ the scene context. Missing factors are marked null rather than guessed. A label definition contains positive visual criteria, exclusion rules, required evidence factors, and an escalation caveat separating visual evidence from a legal conclusion.

The resulting \eventgraph has two edge types. \emph{Hierarchy edges} connect a concept to broader domain and entity nodes. \emph{Confusion edges} connect labels that share most factors but differ in a decision-critical factor, such as floating versus waterside garbage ($r$), building material versus construction waste ($s$), or smoke emission versus suspected fire ($c$ and temporal evidence). These neighbors define hard-negative sampling and diagnostic evaluation. Orthogonally, each unique context induces a \samplegraph
\begin{equation}
  G_i=\left(x_i^g,\{x_{ij}^l\}_{j=1}^{m_i},Y_i,A_i\right),
\end{equation}
where $Y_i$ is the union of valid context events and $A_i\subseteq\{1,\ldots,m_i\}\times Y_i$ links evidence crops to events. This many-to-many graph defines a context-set task in addition to pair recognition.

\subsection{Reconciled Snapshot and Release Tiers}

The business catalog contains 67 concepts and 26,749 accumulated annotation records; those records are not verified unique images. A byte-level filesystem audit finds 9,093 filename-matched context/crop annotations across 45 present concepts, with no orphan or undecodable pair. Exact hashes reduce these rows to 7,279 unique contexts, 9,056 unique crops, and 9,082 unique pairs. The visible distribution remains extreme: the largest concept contributes 73.9\% of pair annotations, the largest two contribute 81.3\%, 41 catalog concepts have fewer than ten visible pairs, and 22 have none. Forty-five class-level download logs expose 1,549 unique business record IDs, but none occurs in the filename or pair manifest, yielding zero deterministic record-to-context joins. \Cref{tab:dataset_audit} reports the versioned audit snapshot.

\begin{table}[t]
\centering
\caption{Content-audited development snapshot. Counts are pair annotations unless marked as exact-hash content groups.}
\label{tab:dataset_audit}
\footnotesize
\begin{tabularx}{\columnwidth}{@{}Y r@{}}
\toprule
Quantity & Value \\
\midrule
Catalog concepts / raw records & 67 / 26,749 \\
Present / absent concepts & 45 / 22 \\
Pair annotations & 9,093 \\
Unique context / crop / pair content & 7,279 / 9,056 / 9,082 \\
Multi-evidence groups / annotations & 1,389 / 3,197 \\
Multi-label groups / annotations & 269 / 623 \\
Redundant exact pairs & 11 (0.12\%) \\
dHash review groups / annotations & 223 / 644 \\
Outside named folder date window & 7,778 (85.5\%) \\
Class-level source IDs / image joins & 1,549 / 0 \\
Verified source media / linked contexts & 0 / 0 \\
Provenance structural / scientific gate & PASS / BLOCKED \\
Proxy components / largest share & 42 / 37.1\% \\
Proxy 30/10/10 support concepts & 4 \\
\bottomrule
\end{tabularx}
\end{table}


We therefore avoid an indefensible 67-way supervised split and define three tiers:
\begin{itemize}
  \item \textbf{Core.} A concept enters supervised evaluation only if it has at least 60 deduplicated source groups, at least three true acquisition groups, and enough groups to place ten or more in each scored split. No concept is finalized until flight/site provenance is recovered. A proxy 30/10/10 context-count diagnostic is reported only as split feasibility, never as Core eligibility.
  \item \textbf{Tail.} Present concepts that fail the core gate are retained for text-only zero-shot, definition-plus-exemplar few-shot, retrieval, and data-acquisition evaluation. They are never oversampled into a conventional test set.
  \item \textbf{Catalog.} Concepts without reconciled or releasable images remain ontology nodes and collection targets, not scored classes. This prevents taxonomy ambition from being misreported as benchmark coverage.
\end{itemize}

\begin{table*}[t]
\centering
\caption{Positioning against representative aerial resources. Scale alone is not our claim; \dataset contributes paired context/evidence, operational factorization, and leakage-aware groups. Counts follow the cited papers and our reconciled development snapshot.}
\label{tab:datasets}
\footnotesize
\begin{tabularx}{\textwidth}{@{}l r c Y c c@{}}
\toprule
Dataset & Scale & Real & Primary scope & Paired context/evidence & Operational factors \\
\midrule
VisDrone~\cite{zhu2018vision} & 10,209 images & Yes & detection and tracking & No & No \\
ERA / CapERA~\cite{mou2020era,papadopoulos2023capera} & 2,864 videos & Mixed & 25 events and captions & No & No \\
UAV-Human~\cite{li2021uavhuman} & 67,428 videos & Yes & human behavior & No & No \\
UAVBench~\cite{zhan2026uavbench} & 261k images & Yes & 10 image/region VLM tasks & Region annotations & No \\
SpatialSky~\cite{zhang2026spatialsky} & 1M samples & -- & 13 UAV spatial categories & Multi-granularity & No \\
UAVLight~\cite{du2026uavlight} & repeat paths & Yes & illumination-robust 3D reconstruction & Multi-time views & No \\
UAVReason~\cite{sun2026uavreason} & 273k samples & No & 22 spatial/temporal reasoning types & Multi-frame & No \\
\dataset (ours) & 9,093 pairs / 7,279 contexts & Yes & operational event sets & Yes & ontology + sample graph \\
\bottomrule
\end{tabularx}
\end{table*}


\subsection{Context--Evidence Sets and Evidence Cards}

For each pair annotation, $x^g$ preserves the wide inspection context and $x^l$ is a linked crop containing putative local evidence. Exact-content grouping reveals 1,389 contexts with multiple evidence crops (3,197 pair annotations) and 269 contexts with multiple labels (623 annotations). Only 11 redundant context--crop pairs exist; thus most repeated contexts encode useful multi-instance or multi-event structure rather than file duplication. One byte-identical pair appears under two labels and is held out pending adjudication; if both labels are valid, it becomes a multi-label target rather than contradictory supervision.

The pair track enables context-only, crop-only, and joint inputs. The set track receives one context and all linked oracle crops, then predicts $Y_i$ and the evidence assignment $A_i$. At deployment, crops may come from an existing detector or retrieval proposal; both tracks separately report oracle- and proposed-crop performance so localization failure is not hidden inside recognition.

We instantiate the remaining uncertainty as four deterministic, row-preserving review ledgers: 623 annotations in 269 multi-label context cases; 644 members in 223 nonexact near-duplicate cases; two annotations in the one exact context--crop cross-label conflict; and one provenance-recovery record for each of the 7,279 exact contexts. Hashes select cases but never decide semantics or provenance. Controlled vocabularies separately encode annotation actions, evidence support, duplicate relations, and provenance status. Two isolated reviewer assignments expose the linked images in an offline visual interface; a strict merger checks identity, row coverage, immutable fields, and controlled values before generating an adjudicator-signed ledger. Agreement is never promoted to ground truth automatically, and excluded rows remain so every frozen manifest can be reproduced.

Every released sample is accompanied by an evidence card with: canonical event and five factors; source, flight, location, and time group identifiers when recoverable; crop-to-source linkage; positive evidence; excluded alternatives; required evidence type (object, region, state, relation, context, or temporal); and review status. A small, stratified grounding subset receives a human-checked box or polygon. Samples whose policy meaning cannot be inferred from a still image are marked \texttt{needs\_temporal\_evidence} and excluded from still-image claims.

\subsection{Caption Construction and Audit}
\label{sec:captiondata}

For a verified pair $(x^g,x^l,y)$, a teacher VLM first produces a label-agnostic description $c^{(0)}$ of visible content. A constrained rewrite then receives $c^{(0)}$, the event definition, exclusions, and verified factors to produce a concise evidence statement $c$. The rewrite must name the event, cite only visible evidence, distinguish the nearest confusion label, and state uncertainty when a required factor is not visible.

Automatic gates reject missing event phrases, factor contradictions, unsupported named entities, and low image--text entailment. A counterfactual audit replaces $y$ with graph neighbor $y^{-}$; a captioner that invents support for $y^{-}$ is flagged. Human review covers every core test item, every rare/tail item used for evaluation, and a stratified sample of training captions. We will report acceptance, correction, and counterfactual-susceptibility rates rather than treating synthetic text as ground truth by default.

\subsection{Leakage-Resistant Splits and Versioning}

Exact hashes, perceptual hashes, and embedding-neighbor search identify duplicates and near duplicates. The exact-context group is the minimum indivisible split unit; all 223 nonexact dHash candidate groups (644 rows) require review and conservatively remain in one partition. Content identity and acquisition identity are separate axes: exact deduplication cannot rule out flight- or site-level shortcuts~\cite{zhang2026leakageaudit}. Final groups are formed in decreasing priority from original media identifier, flight, location, and time window, with every crop and its context together. Entire groups are assigned to train, validation, seen-location test, or unseen-location test. Hyperparameters are frozen on validation data, and the test set is evaluated once per seed.

We make the unresolved gap measurable through a four-level provenance ladder: A requires verified source media, flight, site, time, and permission; B retains verified media and time but incomplete geography; C contains filename-date and content-group proxies only; and D has no parseable date. All current contexts are C or D. For pipeline debugging, a provisional split unions contexts sharing an inferred date or dHash candidate group, producing 42 indivisible components with zero exact, near, or inferred-date cross-split hits. This is still unusable for main results: the largest component holds 2,691 contexts (37.1\%), only four concepts meet a weak 30/10/10 context-support diagnostic, and 47/44 catalog concepts have zero validation/test contexts. These numbers are evidence that filename dates are too coarse, not evidence of leakage-free evaluation.

Revision~6 converts that warning into a fail-closed release gate with three normalized grains: 1,549 business source records, a canonical source-media acquisition registry, and 7,279 exact context-to-media dispositions. The frozen templates pass structural validation, but the scientific gate is \texttt{BLOCKED}: 0/1,549 records are resolved, the media registry has 0 verified rows, and 0/7,279 contexts are linked or excluded. The official split builder therefore emits no manifest. It becomes executable only when every record and context is dispositioned, linked contexts resolve to reviewed media with flight/site/time/permission/privacy evidence, and at least four sites and eight flights exist. A synthetic complete fixture exercises the \texttt{READY} path and asserts zero source-media, flight, site, exact-content, and reviewed-near-group leakage; it is a software test, never dataset evidence.

The release includes a machine-readable data card, taxonomy version, split manifest, file checksums, duplicate-cluster identifiers, adjudication ledgers, and evaluation code. Any later data addition creates a new version; test groups are never silently moved into training. Sensitive coordinates and identifiers are removed or coarsened, and image release is conditional on permission and privacy review.
